\documentclass{article}
\usepackage[letterpaper, margin=1in]{geometry}
\usepackage{amsmath}

\title{UK Derivation}
\author{Cameron Wolfe}
\date{\today}

\begin{document}

\section{UK Constraints Derivation}

We will simulate two six degree of freedom bodies joined by a
revolute joint.  Let the first object be object 1, and the second
object be object 2.  Each object has a position and an attitude.  For
each object, the state is comprised of the inertial position and the
inertial to body orientation represented as MRPs.

\begin{align}
    x_i &=
    \begin{bmatrix}
        r_{i/0}^0 \\
        \sigma_0^i
    \end{bmatrix}
\end{align}

The overall state being the concatenation of these two states:

\begin{align}
    x &=
    \begin{bmatrix}
        x_1 \\
        x_2
    \end{bmatrix} =
    \begin{bmatrix}
        r_{1/0}^0 \\
        \sigma_0^1 \\
        r_{2/0}^0 \\
        \sigma_0^2
    \end{bmatrix}
\end{align}

There are two parts of the constraint, positional and rotational.
Let $r_{j/1}^1$ represent the position of the joint on body one in
the body frame, and $j_1$ be the direction vector of the joint
on body one in the body frame.

The position and rotation constraints are:

\begin{align}
    \phi_r &= r_{1/0}^0 + C((\sigma_0^1)^{-1})r_{j/1}^1 - r_{2/0}^0 -
    C((\sigma_0^2)^{-1})r_{j/2}^2 \\
    \phi_\sigma &=
    (C((\sigma_0^1)^{-1})j_1)\times(C((\sigma_0^2)^{-1})j_2)
\end{align}

Where $C(\sigma)$ is the DCM derived from the MRPs (using the
following formula):

\begin{align}
    C(\sigma) &= I -
    \frac{4(1-\sigma^T\sigma)}{(1+\sigma^T\sigma)^2}\tilde{\sigma} +
    \frac{8}{(1+\sigma^T\sigma)^2}\tilde{\sigma}^2
\end{align}

With these two constraints defined, we can take the second derivative to get the
constraint equations $A\ddot{x}=b$.

\begin{align}
    \phi_r &= r_{1/0}^0 + C((\sigma_0^1)^{-1})r_{j/1}^1 - r_{2/0}^0 -
    C((\sigma_0^2)^{-1})r_{j/2}^2 \\
    \begin{split}
        \frac{d\phi_r}{dt} &= \dot{r}_{1/0}^0 + \dot{C}((\sigma_0^1)^{-1})r_{j/1}^1 -
        \dot{r}_{2/0}^0 - \dot{C}((\sigma_0^2)^{-1})r_{j/2}^2 \\
        &= \dot{r}_{1/0}^0 + C((\sigma_0^1)^{-1})\tilde{\omega}_{1/0}^1r_{j/1}^1 -
        \dot{r}_{2/0}^0 - C((\sigma_0^2)^{-1})\tilde{\omega}_{2/0}^2r_{j/2}^2
    \end{split} \\
    \begin{split}
        \frac{d^2\phi_r}{dt^2} &= \ddot{r}_{1/0}^0 +
        C((\sigma_0^1)^{-1})\tilde{\alpha}_{1/0}^1r_{j/1}^1 +
        C((\sigma_0^1)^{-1})(\tilde{\omega}_{1/0}^1)^2r_{j/1}^1 \\
        &- \ddot{r}_{2/0}^0 - C((\sigma_0^2)^{-1})\tilde{\alpha}_{2/0}^2r_{j/2}^2 -
        C((\sigma_0^2)^{-1})(\tilde{\omega}_{2/0}^2)^2r_{j/2}^2
    \end{split} \\
    \begin{split}
        \label{eqn:ddot_phi_r}
        \frac{d^2\phi_r}{dt^2} &= \ddot{r}_{1/0}^0 -
        C((\sigma_0^1)^{-1})\tilde{r}_{j/1}^1\alpha_{1/0}^1 +
        C((\sigma_0^1)^{-1})(\tilde{\omega}_{1/0}^1)^2r_{j/1}^1 \\
        &- \ddot{r}_{2/0}^0 + C((\sigma_0^2)^{-1})\tilde{r}_{j/2}^2\alpha_{2/0}^2 -
        C((\sigma_0^2)^{-1})(\tilde{\omega}_{2/0}^2)^2r_{j/2}^2
    \end{split}
\end{align}

Before stepping into $\phi_\sigma$, let us define some terms that will be useful, and
take their derivatives:

\begin{align}
    u &= C((\sigma_0^1)^{-1})j_1 \\
    \dot{u} &= C((\sigma_0^1)^{-1})\tilde{\omega}_{1/0}^1j_1 \\
    \begin{split}
        \ddot{u} &= C((\sigma_0^1)^{-1})\tilde{\alpha}_{1/0}^1j_1 +
        C((\sigma_0^1)^{-1})(\tilde{\omega}_{1/0}^1)^2j_1 \\
        &= -C((\sigma_0^1)^{-1})\tilde{j}_1\alpha_{1/0}^1 +
        C((\sigma_0^1)^{-1})(\tilde{\omega}_{1/0}^1)^2j_1
    \end{split} \\
    v &= C((\sigma_0^2)^{-1})j_2 \\
    \dot{v} &= C((\sigma_0^2)^{-1})\tilde{\omega}_{2/0}^2j_2 \\
    \begin{split}
        \ddot{v} &= C((\sigma_0^2)^{-1})\tilde{\alpha}_{2/0}^2j_2 +
        C((\sigma_0^2)^{-1})(\tilde{\omega}_{2/0}^2)^2j_2 \\
        &= -C((\sigma_0^2)^{-1})\tilde{j}_2\alpha_{2/0}^2 +
        C((\sigma_0^2)^{-1})(\tilde{\omega}_{2/0}^2)^2j_2
    \end{split}
\end{align}

\begin{align}
    \phi_\sigma &= u\times v\\
    \begin{split}
        \frac{d\phi_\sigma}{dt} &= \dot{u}\times v + u \times\dot{v}
        \\&=
        (C((\sigma_0^1)^{-1})\tilde{\omega}_{1/0}^1j_1)\times v +
        u\times(C((\sigma_0^2)^{-1})\tilde{\omega}_{2/0}^2j_2)
    \end{split} \\
    \begin{split}
        \label{eqn:ddot_phi_sigma}
        \frac{d^2\phi_\sigma}{dt^2} &= \ddot{u}\times v + 2\dot{u}\times\dot{v} + u\times\ddot{v} \\
        &= (-C((\sigma_0^1)^{-1})\tilde{j}_1\alpha_{1/0}^1 +
        C((\sigma_0^1)^{-1})(\tilde{\omega}_{1/0}^1)^2j_1) \times v \\
        &+
        2(C((\sigma_0^1)^{-1})\tilde{\omega}_{1/0}^1j_1) \times
        (C((\sigma_0^2)^{-1})\tilde{\omega}_{2/0}^2j_2) \\
        &+ u\times(-C((\sigma_0^2)^{-1})\tilde{j}_2\alpha_{2/0}^2 +
        C((\sigma_0^2)^{-1})(\tilde{\omega}_{2/0}^2)^2j_2)
    \end{split}
\end{align}

In their current form, we cannot use these expressions, as we need to get rid of the angular
acceleration terms (and replace them with $\ddot{\sigma})$.  We can start with the
dynamics equation for MRPs relating $\dot{\sigma}$ to $\omega$, and differentiate it.

\begin{align}
    \omega &= 4B^{-1}\dot{\sigma} = 4\frac{(1-\sigma^T\sigma)I - 2\tilde{\sigma} +
    2\sigma\sigma^T}{(1+\sigma^T\sigma)^2}\dot{\sigma} \\
    \begin{split}
        \alpha &= 4B^{-1}\ddot{\sigma} + 4\dot{B}^{-1}\dot{\sigma}
        \label{eqn:alpha_to_sigma}
    \end{split} \\
    \dot{B}^{-1} &= \frac{(-2\sigma^T\dot{\sigma}I - 2\tilde{\dot{\sigma}} +
        2\dot{\sigma}\sigma^T + 2\sigma\dot{\sigma}^T)(1 + \sigma^T\sigma) -
    4\sigma^T\dot{\sigma}((1 - \sigma^T\sigma)I - 2\tilde{\sigma} + 2\sigma\sigma^T)}{(1
    + \sigma^T\sigma)^3}
\end{align}

Plugging Equation \eqref{eqn:alpha_to_sigma} into Equations \eqref{eqn:ddot_phi_r} and
\eqref{eqn:ddot_phi_sigma}, we get:

\begin{align}
    \begin{split}
        \frac{d^2\phi_r}{dt^2} &= \ddot{r}_{1/0}^0 -
        C((\sigma_0^1)^{-1})\tilde{r}_{j/1}^1(4B(\sigma_0^1)^{-1}\ddot{\sigma}_0^1 +
        4(\dot{B}(\sigma_0^1))^{-1}\dot{\sigma}_0^1) +
        C((\sigma_0^1)^{-1})(\tilde{\omega}_{1/0}^1)^2r_{j/1}^1 \\
        &- \ddot{r}_{2/0}^0 +
        C((\sigma_0^2)^{-1})\tilde{r}_{j/2}^2(4B(\sigma_0^2)^{-1}\ddot{\sigma}_0^2 +
        4(\dot{B}(\sigma_0^2))^{-1}\dot{\sigma}_0^2) -
        C((\sigma_0^2)^{-1})(\tilde{\omega}_{2/0}^2)^2r_{j/2}^2
    \end{split} \\
    \begin{split}
        \frac{d^2\phi_\sigma}{dt^2} &=
        (-C((\sigma_0^1)^{-1})\tilde{j}_1(4B(\sigma_0^1)^{-1}\ddot{\sigma}_0^1 +
            4(\dot{B}(\sigma_0^1))^{-1}\dot{\sigma}_0^1) +
        C((\sigma_0^1)^{-1})(\tilde{\omega}_{1/0}^1)^2j_1) \times v \\
        &+ 2(C((\sigma_0^1)^{-1})\tilde{\omega}_{1/0}^1j_1) \times
        (C((\sigma_0^2)^{-1})\tilde{\omega}_{2/0}^2j_2) \\
        &+ u\times(-C((\sigma_0^2)^{-1})\tilde{j}_2(4B(\sigma_0^2)^{-1}\ddot{\sigma}_0^2 +
            4(\dot{B}(\sigma_0^2))^{-1}\dot{\sigma}_0^2) +
        C((\sigma_0^2)^{-1})(\tilde{\omega}_{2/0}^2)^2j_2)
    \end{split}
\end{align}

Using these, we can get $A$ and $b$ for our UK constraint equation

\begin{align}
    A =
    \begin{bmatrix}
        I_3 & -4C((\sigma_0^1)^{-1})\tilde{r}_{j/1}^1B(\sigma_0^1)^{-1} & -I_3 &
        4C((\sigma_0^2)^{-1})\tilde{r}_{j/2}^2B(\sigma_0^2)^{-1} \\
        0_{3\times3} & 4\tilde{v}C((\sigma_0^1)^{-1})\tilde{j}_1B(\sigma_0^1)^{-1} &
        0_{3\times3} & -4\tilde{u}C((\sigma_0^2)^{-1})\tilde{j}_2B(\sigma_0^2)^{-1}
    \end{bmatrix}
\end{align}
\begin{align}
    \begin{split}
        b_r &=
        4C((\sigma_0^1)^{-1})\tilde{r}_{j/1}^1(\dot{B}(\sigma_0^1))^{-1}\dot{\sigma}_0^1
        - C((\sigma_0^1)^{-1})(\tilde{\omega}_{1/0}^1)^2r_{j/1}^1 \\
        &-
        4C((\sigma_0^2)^{-1})\tilde{r}_{j/2}^2(\dot{B}(\sigma_0^2))^{-1}\dot{\sigma}_0^2
        + C((\sigma_0^2)^{-1})(\tilde{\omega}_{2/0}^2)^2r_{j/2}^2
    \end{split} \\
    \begin{split}
        b_\sigma &=
        -4\tilde{v}C((\sigma_0^1)^{-1})\tilde{j}_1(\dot{B}(\sigma_0^1))^{-1}\dot{\sigma}_0^1
        + \tilde{v}C((\sigma_0^1)^{-1})(\tilde{\omega}_{1/0}^1)^2j_1 \\
        &- 2(C((\sigma_0^1)^{-1})\tilde{\omega}_{1/0}^1j_1) \times
        (C((\sigma_0^2)^{-1})\tilde{\omega}_{2/0}^2j_2) \\
        &+
        4\tilde{u}C((\sigma_0^2)^{-1})\tilde{j}_2(\dot{B}(\sigma_0^2))^{-1}\dot{\sigma}_0^2
        - \tilde{u}C((\sigma_0^2)^{-1})(\tilde{\omega}_{2/0}^2)^2j_2
    \end{split}
\end{align}

\section{UK Dynamics Derivation}

The equations of motion for this rigidbody can be broken into two parts, the translational equations of motion, and the rotational equations of motion.  The translational equations are pretty simple.

\begin{align}
    m_i\ddot{r}_{i/0}^0 &= F
\end{align}

The rotational equations of motion come from euler equations:

\begin{align}
    I_i^i\dot{\omega} &= M - \omega\times I_i^i\omega
\end{align}

However, we need to get this in terms of $\ddot{\sigma}_0^i$.

\begin{align}
    \dot{\sigma} &= \frac{B(\sigma)}{4}\omega \\
    \ddot{\sigma} &= \frac{1}{4}(\dot{B}(\sigma)\omega + B\dot{\omega}) \\
                  &= \frac{1}{4}(\dot{B}(\sigma)\omega + B((I_i^i)^{-1}M - (I_i^i)^{-1}\tilde{\omega}I_i^i\omega))
\end{align}

This gives us the overall equations of motion (per body) as

\begin{align}
    \begin{bmatrix}
        m_iI_3 & 0_{3 \times 3} \\
        0_{3 \times 3} & 4I_3
        \end{bmatrix} \begin{bmatrix}
        \ddot{r}_{i/0}^0 \\
        \ddot{\sigma}_0^i
        \end{bmatrix} = \begin{bmatrix}
        F \\
        \dot{B}(\sigma)\omega + B((I_i^i)^{-1}M - (I_i^i)^{-1}\tilde{\omega}I_i^i\omega)
    \end{bmatrix}
\end{align}

\end{document}
